% !TeX root = ../高等代数教材（谢启鸿）笔记.tex
% 由主文档通过 \input 载入；请编译主文档。
\setcounter{chapter}{5} % 保留原书章号。
\chapter{特征值}
\label{chap:xie-eigenvalues}

本章研究有限维线性空间上的线性变换能否通过选择合适的基得到简单的表示矩阵。若一组基中的每个向量都被变换为自身的数乘，表示矩阵便是对角阵；一般情况下，先寻找这样的特殊向量，仍能揭示变换的不变子空间和相似结构。本章依次讨论特征值与特征向量、可对角化的条件、矩阵所满足的多项式恒等式，以及无需精确求根的特征值估计。

除特别说明外，$K$ 表示复数域 $\mathbb C$ 的一个子域，$V$ 是 $K$ 上的 $n$ 维线性空间，$n\geq1$，$T:V\to V$ 为线性变换，$I_V$ 为恒等变换，$I_n$ 为 $n$ 阶单位矩阵。矩阵的特征向量均采用列向量；讨论全部复特征值时按特征多项式的根的重数计数。

\section{特征值和特征向量}
\label{sec:xie-eigenvectors}
% 依据教材第 267--274 页正文及习题 6.1。

若非零向量 $v$ 张成的一维子空间在 $T$ 下保持不变，则必有 $Tv=\lambda v$，其中 $\lambda\in K$。因此，寻找一维不变子空间等价于寻找变换只作数乘的方向；这里的“方向”是代数意义上的一维子空间，并不限于实数空间中的几何方向。

\begin{definition}
\label{def:xie-eigenvalue}
若存在 $\lambda\in K$ 和 $v\in V\setminus\{0\}$，使 $Tv=\lambda v$，则称 $\lambda$ 为 $T$ 的\emph{特征值}（eigenvalue），$v$ 为属于 $\lambda$ 的\emph{特征向量}（eigenvector）。对应的\emph{特征子空间}（eigenspace）定义为
\[
 E_\lambda(T)=\ker(T-\lambda I_V)=\{v\in V:Tv=\lambda v\}.
\]
对 $A\in M_n(K)$，相应定义为 $Av=\lambda v$，其中 $v\in K^n\setminus\{0\}$，并记 $E_\lambda(A)=\ker(A-\lambda I_n)$。
\end{definition}

核是子空间，且 $T$ 在 $E_\lambda(T)$ 上的限制就是 $\lambda I$，所以特征子空间也是 $T$ 的不变子空间。\emph[caution]{零向量属于每个特征子空间，但不是特征向量}；非零条件使“存在特征向量”具有内容。若 $A$ 是 $T$ 在某组基下的表示矩阵，$x$ 是 $v$ 的坐标列，则 $Tv=\lambda v$ 等价于 $Ax=\lambda x$，因此两种表述完全对应。

\begin{definition}
\label{def:xie-characteristic-polynomial}
矩阵 $A$ 的\emph{特征多项式}（characteristic polynomial）定义为
\begin{equation}
\label{eq:xie-characteristic-polynomial}
 \chi_A(t)=\det(tI_n-A).
\end{equation}
它是 $K[t]$ 中的 $n$ 次首一多项式，$tI_n-A$ 称为特征矩阵。
\end{definition}

齐次方程 $(\lambda I_n-A)x=0$ 有非零解，当且仅当其系数矩阵不可逆。因此，\emph{$\lambda\in K$ 是 $A$ 的特征值，当且仅当 $\chi_A(\lambda)=0$}。计算时先求特征多项式在指定数域中的根，再对每个不同的根求齐次方程的基础解系；该基础解系张成特征子空间，去掉零向量即得到全部特征向量。

\begin{theorem}[相似不变性]
\label{thm:xie-characteristic-similarity}
若 $B=P^{-1}AP$，其中 $P\in M_n(K)$ 可逆，则 $\chi_B(t)=\chi_A(t)$，且对任意 $\lambda\in K$，有
\[
 E_\lambda(B)=P^{-1}E_\lambda(A).
\]
特别地，相似矩阵有相同的特征值、代数重数和特征子空间维数。
\end{theorem}
\begin{proof}
利用行列式的乘法性质，
\[
 \det(tI_n-B)=\det\bigl(P^{-1}(tI_n-A)P\bigr)=\det(tI_n-A).
\]
又有 $(B-\lambda I_n)x=0$ 当且仅当 $(A-\lambda I_n)Px=0$，由 $P$ 可逆即得子空间关系。
\end{proof}

同一个线性变换在不同基下的表示矩阵相似，故可用任意表示矩阵定义 $T$ 的特征多项式 $\chi_T(t)$；这个定义与基的选择无关。设 $A=(a_{ij})$，记 $\operatorname{tr}A=\sum_i a_{ii}$ 为矩阵的迹（trace）。当 $n\geq2$ 时，展开行列式可得
\[
 \chi_A(t)=t^n-(\operatorname{tr}A)t^{n-1}+\cdots+(-1)^n\det A.
\]
$n=1$ 时，$\chi_A(t)=t-a_{11}$。在 $\mathbb C$ 中将特征多项式分解为 $\prod_{i=1}^n(t-\lambda_i)$，比较系数，对所有 $n\geq1$ 都有
\begin{equation}
\label{eq:xie-eigen-trace-det}
 \sum_{i=1}^n\lambda_i=\operatorname{tr}A,
 \qquad \prod_{i=1}^n\lambda_i=\det A.
\end{equation}
这里的根按代数重数重复列出。即使 $A$ 的元素在 $K$ 中，其全部根也未必都在 $K$ 中；和、积公式可在 $\mathbb C$ 中使用。

\begin{example}
\label{ex:xie-eigen-triangular}
求上三角矩阵 $A=(a_{ij})\in M_n(K)$ 的特征值。
\end{example}
\begin{solve}
$tI_n-A$ 仍为上三角矩阵，故 $\chi_A(t)=\prod_{i=1}^n(t-a_{ii})$。因此全部特征值就是主对角线元素，重复出现的次数等于对应的代数重数。下三角矩阵有相同结论。
\end{solve}

\begin{example}
\label{ex:xie-eigen-defective-three}
求矩阵
\[
 A=\begin{pmatrix}3&1&-1\\2&2&-1\\2&2&0\end{pmatrix}
\]
的特征值与特征子空间。
\end{example}
\begin{solve}
直接计算得 $\chi_A(t)=t^3-5t^2+8t-4=(t-1)(t-2)^2$。分别求解 $(A-I_3)x=0$ 与 $(A-2I_3)x=0$，得到
\[
 E_1(A)=\operatorname{span}\{(1,0,2)^{\mathsf T}\},
 \qquad E_2(A)=\operatorname{span}\{(1,1,2)^{\mathsf T}\}.
\]
例如第一个方程组化简为 $x_2=0$、$x_3=2x_1$；第二个化简为 $x_2=x_1$、$x_3=2x_1$。各子空间中的非零向量就是对应的全部特征向量。特征值 $2$ 是二重根，但对应的特征子空间只有一维，这预示着重根未必提供同样多的线性无关特征向量。
\end{solve}

\begin{example}
\label{ex:xie-eigen-field}
考察 $R=\begin{pmatrix}0&-1\\1&0\end{pmatrix}$ 在不同数域上的特征值。
\end{example}
\begin{solve}
$\chi_R(t)=t^2+1$，所以 $R$ 在 $\mathbb R$ 上没有特征值，在 $\mathbb C$ 上的特征值为 $i,-i$，对应特征子空间分别由 $(1,-i)^{\mathsf T}$ 和 $(1,i)^{\mathsf T}$ 张成。因而讨论特征值或对角化时必须指定数域。
\end{solve}

\begin{theorem}[上三角化]
\label{thm:xie-triangularization}
$A\in M_n(K)$ 在 $K$ 上可作\emph{上三角化}（triangularization），即相似于上三角矩阵，当且仅当 $\chi_A(t)$ 在 $K$ 上分裂为一次因式的乘积。特别地，每个复方阵都相似于上三角矩阵。
\end{theorem}
\begin{proof}
必要性由例~\ref{ex:xie-eigen-triangular} 和相似不变性得到。证明充分性，对 $n$ 归纳；$n=1$ 时成立。取一个特征值 $\lambda_1\in K$ 及其非零特征向量 $v_1$，将 $v_1$ 扩充为 $K^n$ 的一组基，组成可逆矩阵 $P$，则
\[
 P^{-1}AP=\begin{pmatrix}\lambda_1&b\\0&A_1\end{pmatrix},
 \qquad \chi_A(t)=(t-\lambda_1)\chi_{A_1}(t).
\]
由分裂条件，$\chi_{A_1}$ 也在 $K$ 上分裂。由归纳假设可取 $Q\in M_{n-1}(K)$ 使 $Q^{-1}A_1Q$ 为上三角矩阵，再用 $\operatorname{diag}(1,Q)$ 作相似变换便得到所需形式。复数域上的分裂性由代数基本定理保证。
\end{proof}

\begin{proposition}[多项式作用于特征值]
\label{prop:xie-polynomial-spectrum}
设 $A\in M_n(K)$ 的全部复特征值为 $\lambda_1,\ldots,\lambda_n$，$p\in K[t]$。则 $p(A)$ 的全部复特征值为 $p(\lambda_1),\ldots,p(\lambda_n)$，按代数重数计数；若 $Av=\lambda v$，则 $p(A)v=p(\lambda)v$。因此，$p(A)=0$ 蕴含每个特征值都满足 $p(\lambda)=0$；$A$ 可逆当且仅当零不是其复特征值，可逆时 $A^{-1}$ 的全部特征值为 $\lambda_1^{-1},\ldots,\lambda_n^{-1}$。
\end{proposition}
\begin{proof}
由 $A^jv=\lambda^jv$ 对 $j$ 归纳，得到关于特征向量的等式；$v\ne0$ 保证 $p(A)=0$ 时 $p(\lambda)=0$。为确定全部特征值及其重数，在 $\mathbb C$ 上取 $P$ 使 $U=P^{-1}AP$ 为上三角矩阵，其对角元为 $\lambda_1,\ldots,\lambda_n$。由
\[
 P^{-1}p(A)P=p(U)
\]
及上三角矩阵乘法的对角元规则，$p(U)$ 的对角元为 $p(\lambda_i)$，故结论成立。可逆性的判据由式~\eqref{eq:xie-eigen-trace-det} 得到；可逆时 $U^{-1}$ 的对角元为 $\lambda_i^{-1}$，故其逆矩阵的结论也成立。
\end{proof}


\subsection{习题 6.1 解答}
\label{subsec:xie-exercises-61}

以下保留教材习题题号，将题面与解答放在一起。涉及非实特征值的题目在 $\mathbb C$ 上求解；若限定在其他数域，则只保留属于该数域的特征值及相应向量。

\begin{example}
\label{ex:xie-61-01}
教材习题 6.1 第 1 题：求下列矩阵的特征值与特征向量：
\[
 A_1=\begin{pmatrix}3&2&2\\2&3&2\\-4&-4&-3\end{pmatrix},\quad
 A_2=\begin{pmatrix}2&-1&2\\5&-3&3\\-1&0&-2\end{pmatrix},\quad
 A_3=\begin{pmatrix}3&-1&-1\\7&-2&-3\\3&-1&-1\end{pmatrix}.
\]
\end{example}
\begin{solve}
计算特征多项式并行化简各特征矩阵，分别得到
\[
 \begin{aligned}
 \chi_{A_1}(t)&=(t-1)^3,& E_1(A_1)&=\{x:x_1+x_2+x_3=0\}
 =\operatorname{span}\{(-1,1,0)^{\mathsf T},(-1,0,1)^{\mathsf T}\},\\
 \chi_{A_2}(t)&=(t+1)^3,& E_{-1}(A_2)&=\operatorname{span}\{(1,1,-1)^{\mathsf T}\},\\
 \chi_{A_3}(t)&=t^3,& E_0(A_3)&=\operatorname{span}\{(1,2,1)^{\mathsf T}\}.
 \end{aligned}
\]
第二个方程组给出 $x_2=x_1$、$x_3=-x_1$，第三个给出 $x_2=2x_1$、$x_3=x_1$，所以没有遗漏其他自由变量。每个所列子空间去掉零向量，就是该矩阵的全部特征向量。
\end{solve}

\begin{example}
\label{ex:xie-61-02}
教材习题 6.1 第 2 题：设 $V=V_1\oplus\cdots\oplus V_s$，各 $V_i$ 都是 $T$ 的不变子空间，$T|_{V_i}$ 的特征多项式为 $f_i(t)$。证明 $\chi_T=\prod_i f_i$；并证明对任一特征值 $\lambda$，有
\[
 E_\lambda(T)=\bigoplus_{i=1}^s\bigl(E_\lambda(T)\cap V_i\bigr).
\]
\end{example}
\begin{solve}
分别取各 $V_i$ 的基并拼接，$T$ 的表示矩阵是由各限制变换的表示矩阵构成的分块对角矩阵，所以特征行列式等于各块特征行列式之积。再将 $v\in V$ 唯一写成 $v=\sum_i v_i$，其中 $v_i\in V_i$。由不变性，$Tv-\lambda v=\sum_i(Tv_i-\lambda v_i)$ 是按这些子空间作出的直和分解，故它为零当且仅当每项为零。这正是所需的特征子空间分解。
\end{solve}

\begin{example}
\label{ex:xie-61-03}
教材习题 6.1 第 3 题：设 $A=\operatorname{diag}(A_1,\ldots,A_s)$，其中 $A_i$ 为 $n_i$ 阶矩阵。证明 $A_i$ 的特征向量可通过在其前后补零延拓为 $A$ 的特征向量；并证明 $E_\lambda(A)$ 的一组基可由各 $E_\lambda(A_i)$ 的基延拓后合并得到。
\end{example}
\begin{solve}
令 $\iota_i:K^{n_i}\to K^{\sum_jn_j}$ 为在第 $i$ 个坐标块放入原向量、其余坐标块补零的映射，则 $A\iota_i(v)=\iota_i(A_iv)$。因此 $A_iv=\lambda v$ 蕴含 $A\iota_i(v)=\lambda\iota_i(v)$，且非零向量延拓后仍非零。一般向量按块写成 $(v_1^{\mathsf T},\ldots,v_s^{\mathsf T})^{\mathsf T}$，特征方程等价于每个 $(A_i-\lambda I_{n_i})v_i=0$，所以
\[
 E_\lambda(A)=\bigoplus_{i=1}^s\iota_i(E_\lambda(A_i)).
\]
非特征值对应的核为零空间，不贡献基向量；其余各块的基延拓后既生成全核，又因坐标块互不重叠而线性无关。
\end{solve}

\begin{example}
\label{ex:xie-61-04}
教材习题 6.1 第 4 题：利用分块对角结构，求下列矩阵的特征值与特征向量：
\[
 B_1=\begin{pmatrix}1&0&0&0\\0&2&0&0\\0&0&a&b\\0&0&c&d\end{pmatrix},\quad
 B_2=\begin{pmatrix}0&1&0&0\\1&0&0&0\\0&0&0&-1\\0&0&1&0\end{pmatrix},\quad
 B_3=\begin{pmatrix}1&3&0&0\\4&2&0&0\\0&0&3&1\\0&0&6&4\end{pmatrix}.
\]
\end{example}
\begin{solve}
对 $B_1$，令 $C=\begin{pmatrix}a&b\\c&d\end{pmatrix}$，则
\[
 \chi_{B_1}(t)=(t-1)(t-2)\bigl((t-a)(t-d)-bc\bigr).
\]
所以特征值为 $1,2$ 及二次式的根 $\frac{a+d\pm\sqrt{(a-d)^2+4bc}}2$，按重数计数。对每个不同的根 $\lambda$，令 $W_\lambda=\ker(C-\lambda I_2)$，则完整的答案为
\[
 E_\lambda(B_1)=\operatorname{span}\{e_j:j\in\{1,2\},\ j=\lambda\}
 \oplus\{(0,0,u,v)^{\mathsf T}:(u,v)^{\mathsf T}\in W_\lambda\}.
\]
其中 $e_j$ 是四维标准单位列向量。若 $C=\lambda I_2$，则 $W_\lambda=K^2$；若 $C-\lambda I_2$ 非零且行列式为零，取其任一非零行 $(r_1,r_2)$，便有 $W_\lambda=\operatorname{span}\{(-r_2,r_1)^{\mathsf T}\}$；若该矩阵可逆，则 $W_\lambda=\{0\}$。这个写法同时处理二次式重根以及它的根与 $1,2$ 重合的情形，不需除以可能为零的参数。

$B_2$ 的特征值为 $1,-1,i,-i$，对应特征子空间依次由
\[
 (1,1,0,0)^{\mathsf T},\quad(1,-1,0,0)^{\mathsf T},\quad
 (0,0,1,-i)^{\mathsf T},\quad(0,0,1,i)^{\mathsf T}
\]
张成。$B_3$ 的两个二阶块的特征多项式分别为 $(t-5)(t+2)$、$(t-1)(t-6)$，故它的特征值为 $5,-2,1,6$，对应特征子空间依次由
\[
 (3,4,0,0)^{\mathsf T},\quad(1,-1,0,0)^{\mathsf T},\quad
 (0,0,1,-2)^{\mathsf T},\quad(0,0,1,3)^{\mathsf T}
\]
张成。各特征子空间中的非零向量即为全部特征向量。
\end{solve}

\begin{example}
\label{ex:xie-61-05}
教材习题 6.1 第 5 题：若 $Av_1=\lambda_1v_1$、$Av_2=\lambda_2v_2$，其中 $v_1,v_2\ne0$ 且 $\lambda_1\ne\lambda_2$，证明 $v_1+v_2$ 不是特征向量。
\end{example}
\begin{solve}
若二者线性相关，则 $v_2=cv_1$（$c\ne0$）会推出 $\lambda_1v_2=Av_2=\lambda_2v_2$，与特征值互异矛盾。因此 $v_1,v_2$ 线性无关，其和非零。若 $A(v_1+v_2)=\mu(v_1+v_2)$，则 $(\lambda_1-\mu)v_1+(\lambda_2-\mu)v_2=0$，线性无关性推出 $\lambda_1=\mu=\lambda_2$，矛盾。
\end{solve}

\begin{example}
\label{ex:xie-61-06}
教材习题 6.1 第 6 题：证明 $n$ 阶矩阵以每个非零列向量为特征向量，当且仅当它是纯量阵。
\end{example}
\begin{solve}
纯量阵 $cI_n$ 满足 $Av=cv$，故充分性成立。反之，标准基向量给出 $Ae_i=c_ie_i$；若 $n\geq2$，对任意 $i\ne j$，向量 $e_i+e_j$ 也为特征向量，于是 $c_ie_i+c_je_j=c_{ij}(e_i+e_j)$，比较系数得 $c_i=c_j$。故所有 $c_i$ 相同，$A=cI_n$。$n=1$ 时结论直接成立。
\end{solve}

\begin{example}
\label{ex:xie-61-07}
教材习题 6.1 第 7 题：设 $A=\begin{pmatrix}2&1&1\\1&2&1\\1&1&a\end{pmatrix}$ 可逆，其伴随矩阵 $\operatorname{adj}A$ 有特征值 $\lambda_0$ 和对应特征向量 $(1,b,1)^{\mathsf T}$，求 $a,b,\lambda_0$。
\end{example}
\begin{solve}
由 $\operatorname{adj}A=(\det A)A^{-1}$，伴随矩阵也可逆，故 $\lambda_0\ne0$，且同一向量是 $A$ 关于 $\mu=\det A/\lambda_0$ 的特征向量。逐行写出方程为
\[
 3+b=\mu,\qquad 2+2b=\mu b,\qquad 1+b+a=\mu.
\]
第一、三式给出 $a=2$，前两式给出 $(b-1)(b+2)=0$。此时 $\det A=4$；$b=1$ 时 $\mu=4$、$\lambda_0=1$，$b=-2$ 时 $\mu=1$、$\lambda_0=4$。所以答案为 $(a,b,\lambda_0)=(2,1,1)$ 或 $(2,-2,4)$，均满足可逆条件。
\end{solve}

\begin{example}
\label{ex:xie-61-08}
教材习题 6.1 第 8 题：设 $A=\begin{pmatrix}a&-1&c\\5&b&3\\1-c&0&-a\end{pmatrix}$，$\det A=-1$，且 $\operatorname{adj}A$ 有特征值 $\lambda_0$ 和对应特征向量 $(-1,-1,1)^{\mathsf T}$，求各参数。
\end{example}
\begin{solve}
同上一题，令 $\mu=\det A/\lambda_0$，则
\[
 1+c-a=-\mu,\qquad -2-b=-\mu,\qquad c-a-1=\mu.
\]
首、末两式相加得 $c=a$，继而 $\mu=-1$、$b=-3$。代入行列式得 $\det A=a-3=-1$，故 $a=c=2$，最后 $\lambda_0=(-1)/(-1)=1$。代回三个方程与行列式条件均成立。
\end{solve}

\begin{example}
\label{ex:xie-61-09}
教材习题 6.1 第 9 题：证明幂零矩阵（nilpotent matrix，即存在正整数 $k$ 使 $A^k=0$）的特征值全为零；对合矩阵（involutory matrix，即 $A^2=I_n$）的特征值只可能为 $1,-1$；幂等矩阵（idempotent matrix，即 $A^2=A$）的特征值只可能为 $0,1$。
\end{example}
\begin{solve}
若 $Av=\lambda v$、$v\ne0$，则 $A^kv=\lambda^kv$。三种恒等式分别推出 $\lambda^k=0$、$\lambda^2=1$、$\lambda^2=\lambda$，从而得到结论。这里“只可能”给出候选集合，不要求集合中的每个数都实际出现。
\end{solve}

\begin{example}
\label{ex:xie-61-10}
教材习题 6.1 第 10 题：设 $\alpha,\beta$ 是非零 $n$ 维列向量，且 $\alpha^{\mathsf T}\beta=0$，求 $A=\alpha\beta^{\mathsf T}$ 的全部特征值。
\end{example}
\begin{solve}
这里转置不含复共轭，故 $\beta^{\mathsf T}\alpha=\alpha^{\mathsf T}\beta=0$。于是 $A^2=\alpha(\beta^{\mathsf T}\alpha)\beta^{\mathsf T}=0$，由上一题知全部 $n$ 个复特征值均为零。一维时两个非零数的乘积不为零，因此题设条件不会发生。
\end{solve}

\begin{example}
\label{ex:xie-61-11}
教材习题 6.1 第 11 题：设 $(A+I_n)^m=0$，其中 $m$ 为正整数，证明 $A$ 可逆并求 $\det A$。
\end{example}
\begin{solve}
任一复特征值满足 $(\lambda+1)^m=0$，故全为 $-1$。于是 $\det A=(-1)^n\ne0$，所以 $A$ 可逆。也可令 $N=A+I_n$，利用有限等比和写出 $A^{-1}=-(I_n+N+\cdots+N^{m-1})$，直接乘回验证。
\end{solve}

\begin{example}
\label{ex:xie-61-12}
教材习题 6.1 第 12 题：设 $A\in M_{m\times n}(K)$、$B\in M_{n\times m}(K)$，$m\geq n$，证明
\begin{equation}
\label{eq:xie-ab-ba-characteristic}
 \det(tI_m-AB)=t^{m-n}\det(tI_n-BA).
\end{equation}
特别地，等阶矩阵的乘积 $AB$ 与 $BA$ 有相同特征多项式。
\end{example}
\begin{solve}
当 $t\ne0$ 时，对分块矩阵 $M=\begin{pmatrix}tI_m&A\\B&I_n\end{pmatrix}$ 分别消去右上块和左下块，得到
\[
 \det M=\det(tI_m-AB)
 =t^m\det(I_n-t^{-1}BA)
 =t^{m-n}\det(tI_n-BA).
\]
这些消去由行列式为 $1$ 的分块单位三角矩阵实现，不改变行列式。两端均为 $t$ 的多项式，在无限数域 $K$ 中对所有非零 $t$ 相等，故为多项式恒等式，$t=0$ 时也成立。
\end{solve}

\begin{example}
\label{ex:xie-61-13}
教材习题 6.1 第 13 题：设 $\alpha,\beta$ 为非零 $n$ 维列向量，求 $I_n-\alpha\beta^{\mathsf T}$ 的全部特征值。
\end{example}
\begin{solve}
令 $s=\beta^{\mathsf T}\alpha$，以 $\alpha$ 为第一个基向量扩充一组基。$\alpha\beta^{\mathsf T}$ 的像包含在 $K\alpha$ 中，且它将 $\alpha$ 映为 $s\alpha$，所以在这组基下是对角元为 $s,0,\ldots,0$ 的上三角矩阵。故所求特征值按重数为 $1-s$ 及 $n-1$ 个 $1$，其特征多项式为 $(t-1+s)(t-1)^{n-1}$；当 $s=0$ 时全部为 $1$，并不能据此断定矩阵就是单位矩阵。
\end{solve}

\begin{example}
\label{ex:xie-61-14}
教材习题 6.1 第 14 题：证明 $AB+B$ 与 $BA+B$ 有相同的特征值。
\end{example}
\begin{solve}
令 $C=A+I_n$，则 $AB+B=CB$、$BA+B=BC$。对 $C,B$ 使用式~\eqref{eq:xie-ab-ba-characteristic} 的等阶情形，得到二者有相同特征多项式，因此特征值连同代数重数都相同。
\end{solve}

\begin{example}
\label{ex:xie-61-15}
教材习题 6.1 第 15 题：若 $P$ 可逆，证明 $B=PAP^{-1}-P^{-1}AP$ 的特征值之和为零。
\end{example}
\begin{solve}
两个被减矩阵都与 $A$ 相似，所以由相似不变性和式~\eqref{eq:xie-eigen-trace-det}，二者的迹都等于 $\operatorname{tr}A$。迹的线性性给出 $\operatorname{tr}B=0$，这就是 $B$ 的全部复特征值按重数计算的和。
\end{solve}

\begin{example}
\label{ex:xie-61-16}
教材习题 6.1 第 16 题：设 $A,B\in M_n(\mathbb C)$ 且 $AB=BA$。证明 $A$ 的各特征子空间在 $B$ 下不变，反之亦然；二者有公共特征向量；并且二者可以同时上三角化。
\end{example}
\begin{solve}
若 $Av=\lambda v$，则 $A(Bv)=B(Av)=\lambda Bv$，所以 $Bv\in E_\lambda(A)$，其中允许 $Bv=0$。交换 $A,B$ 即得另一方向。取 $A$ 的任一非零特征子空间，$B$ 在其上的限制是一个非零维复线性空间上的线性变换，故有非零特征向量；它同时是 $A$ 的特征向量。

最后对 $n$ 归纳。以一个公共特征向量为首向量扩充一组基，在同一个可逆矩阵 $P$ 下有
\[
 P^{-1}AP=\begin{pmatrix}\lambda&a\\0&A_1\end{pmatrix},\qquad
 P^{-1}BP=\begin{pmatrix}\mu&b\\0&B_1\end{pmatrix}.
\]
比较乘积的右下块得 $A_1B_1=B_1A_1$。由归纳假设，存在同一个 $Q$ 使这两个低阶块同时上三角化，再用 $P\operatorname{diag}(1,Q)$ 作变换即可。$n=1$ 为归纳起点。复数域保证每一步限制变换都存在特征值。
\end{solve}

\section{对角化}
\label{sec:xie-diagonalization}
% 依据教材第 276--282 页正文及习题 6.2。

\begin{definition}
\label{def:xie-diagonalizable}
若存在 $V$ 的一组基使 $T$ 的表示矩阵为对角矩阵，则称 $T$ 在 $K$ 上\emph{可对角化}（diagonalizable）。等价地，矩阵 $A\in M_n(K)$ 可对角化是指存在 $P\in M_n(K)$ 可逆，使 $P^{-1}AP$ 为对角矩阵。
\end{definition}

\begin{theorem}[特征向量基判据]
\label{thm:xie-eigenbasis}
$T$ 可对角化，当且仅当 $V$ 存在一组由 $T$ 的特征向量组成的基。相应地，$A$ 可对角化，当且仅当它有 $n$ 个线性无关的特征向量。若 $Av_j=\lambda_jv_j$，则
\begin{equation}
\label{eq:xie-diagonalizer}
 P=(v_1,\ldots,v_n)\text{ 可逆}
 \quad\Longrightarrow\quad
 P^{-1}AP=\operatorname{diag}(\lambda_1,\ldots,\lambda_n).
\end{equation}
\end{theorem}
\begin{proof}
设 $P=(v_1,\ldots,v_n)$，则 $AP=P\operatorname{diag}(\lambda_1,\ldots,\lambda_n)$ 恰好逐列表示 $Av_j=\lambda_jv_j$。$P$ 可逆等价于各列构成一组基，因而上述等式同时证明两个方向。
\end{proof}

在特征向量基 $v_1,\ldots,v_n$ 下，对任意 $v=\sum_i c_iv_i$，有 $Tv=\sum_i\lambda_ic_iv_i$。若把非零特征值对应的基向量排在前面，设 $\lambda_1,\ldots,\lambda_r\ne0$，其余为零，则逐项比较系数可得 $\operatorname{im}T=\operatorname{span}\{v_1,\ldots,v_r\}$、$\ker T=\operatorname{span}\{v_{r+1},\ldots,v_n\}$，从而 $\operatorname{rank}T=r$。可对角化时，像、核和反复作用的计算都化为对各基向量的标量计算。

\begin{example}
\label{ex:xie-jordan-two}
矩阵 $J=\begin{pmatrix}1&1\\0&1\end{pmatrix}$ 是否可对角化？
\end{example}
\begin{solve}
$\chi_J(t)=(t-1)^2$，但 $E_1(J)=\operatorname{span}\{(1,0)^{\mathsf T}\}$，只有一个线性无关特征向量，因此不可对角化。也可注意：若它能对角化，所得对角阵只能是 $I_2$，而 $P^{-1}JP=I_2$ 会推出 $J=I_2$，与题设矛盾。
\end{solve}

\begin{theorem}[不同特征子空间之和为直和]
\label{thm:xie-eigenspaces-direct}
设 $\lambda_1,\ldots,\lambda_s\in K$ 是 $T$ 的互不相同的特征值，则
\[
 E_{\lambda_1}(T)+\cdots+E_{\lambda_s}(T)
 =E_{\lambda_1}(T)\oplus\cdots\oplus E_{\lambda_s}(T).
\]
特别地，从不同特征值的特征子空间中各取一个非零向量，所得向量必线性无关。
\end{theorem}
\begin{proof}
对 $s$ 归纳。$s=1$ 时结论成立。设 $v_i\in E_{\lambda_i}(T)$ 且 $\sum_{i=1}^s v_i=0$，对等式施加 $T-\lambda_sI_V$，得到
\[
 \sum_{i=1}^{s-1}(\lambda_i-\lambda_s)v_i=0.
\]
归纳假设说明每项 $(\lambda_i-\lambda_s)v_i$ 都为零；由特征值互异，$v_1=\cdots=v_{s-1}=0$，进而 $v_s=0$，故和为直和。将 $v_i$ 取为待检验线性组合中的各项，就得到最后的结论。
\end{proof}

\begin{corollary}
\label{cor:xie-distinct-eigenvalues}
若 $\chi_A(t)$ 在 $K$ 上有 $n$ 个互不相同的根，则 $A$ 在 $K$ 上可对角化。一般地，$T$ 可对角化，当且仅当 $V$ 为它在 $K$ 上各个特征子空间的直和。
\end{corollary}
\begin{proof}
第一个结论由不同特征值的特征向量线性无关得到。若各特征子空间的直和是 $V$，分别取各子空间的基，拼接后就是特征向量基；反之，若已有特征向量基，按特征值分组便知这些子空间之和包含一组基，从而等于 $V$，再由定理~\ref{thm:xie-eigenspaces-direct} 知其为直和。
\end{proof}

\begin{definition}
\label{def:xie-multiplicities}
设 $\lambda$ 是 $A$ 的特征值。它作为 $\chi_A(t)$ 的根的重数称为\emph{代数重数}（algebraic multiplicity），记为 $a_\lambda$；特征子空间维数
\[
 g_\lambda=\dim E_\lambda(A)=n-\operatorname{rank}(A-\lambda I_n)
\]
称为\emph{几何重数}（geometric multiplicity），教材也称为该特征值的度数。线性变换的两种重数由其表示矩阵等价定义。
\end{definition}

\begin{lemma}
\label{lem:xie-geometric-algebraic}
每个特征值满足 $1\leq g_\lambda\leq a_\lambda$。
\end{lemma}
\begin{proof}
特征子空间含有非零向量，故 $g_\lambda\geq1$。取它的一组基并扩充为全空间的基，$A$ 在这组基下相似于
\[
 \begin{pmatrix}\lambda I_{g_\lambda}&B\\0&C\end{pmatrix}.
\]
因此 $\chi_A(t)=(t-\lambda)^{g_\lambda}\det(tI-C)$，可见根 $\lambda$ 的重数至少为 $g_\lambda$。若 $g_\lambda=n$，上式直接理解为 $\chi_A(t)=(t-\lambda)^n$。
\end{proof}

\begin{theorem}[重数判据]
\label{thm:xie-multiplicity-criterion}
矩阵 $A\in M_n(K)$ 在 $K$ 上可对角化，当且仅当 $\chi_A(t)$ 在 $K$ 上完全分裂，且它的每个特征值都满足 $g_\lambda=a_\lambda$。
\end{theorem}
\begin{proof}
若 $A$ 可对角化，对角阵的特征多项式在 $K$ 上完全分裂。反之及重数条件的必要性，可统一由维数比较说明：分裂时，设全部不同的特征值为 $\lambda_1,\ldots,\lambda_s$，则
\[
 \dim\bigoplus_{i=1}^sE_{\lambda_i}(A)=\sum_{i=1}^s g_{\lambda_i}
 \leq\sum_{i=1}^s a_{\lambda_i}=n.
\]
由推论~\ref{cor:xie-distinct-eigenvalues}，可对角化等价于最左端等于 $n$；结合逐项不等式，等号成立等价于每项 $g_{\lambda_i}=a_{\lambda_i}$。
\end{proof}

教材用“有完全的特征向量系”描述各特征值的度数等于重数。在复数域上，特征多项式自动完全分裂，这一条件便与可对角化等价；在一般数域上，必须同时保留分裂条件，否则例~\ref{ex:xie-eigen-field} 中的实矩阵没有实特征值，仅检验已有实特征值的重数并不能判定可对角化。有 $n$ 个不同特征值只是充分条件，纯量阵 $cI_n$ 表明它并非必要条件。

\begin{example}
\label{ex:xie-diagonalization-three}
判断 $A=\begin{pmatrix}1&0&0\\-2&5&-2\\-2&4&-1\end{pmatrix}$ 是否可对角化，并求相应的可逆矩阵。
\end{example}
\begin{solve}
$\chi_A(t)=(t-1)^2(t-3)$，直接解方程得到
\[
 E_1(A)=\operatorname{span}\{(2,1,0)^{\mathsf T},(-1,0,1)^{\mathsf T}\},
 \quad E_3(A)=\operatorname{span}\{(0,1,1)^{\mathsf T}\}.
\]
两种特征值的几何重数分别为 $2,1$，均等于代数重数。因而可取
\[
 P=\begin{pmatrix}2&-1&0\\1&0&1\\0&1&1\end{pmatrix},
 \qquad P^{-1}AP=\operatorname{diag}(1,1,3).
\]
构造 $P$ 时，同一特征子空间内可以更换基，各列也可以重新排列，但对角阵的对角元必须与列向量的特征值依次对应。
\end{solve}

\begin{example}
\label{ex:xie-diagonal-power}
用对角化计算 $A=\begin{pmatrix}1&0\\1&-2\end{pmatrix}$ 的 $10$ 次幂。
\end{example}
\begin{solve}
特征值为 $1,-2$，分别取特征向量 $(3,1)^{\mathsf T}$、$(0,1)^{\mathsf T}$，则
\[
 P=\begin{pmatrix}3&0\\1&1\end{pmatrix},\quad
 P^{-1}AP=\operatorname{diag}(1,-2),\quad
 A^m=P\operatorname{diag}(1,(-2)^m)P^{-1}
 =\begin{pmatrix}1&0\\\dfrac{1-(-2)^m}{3}&(-2)^m\end{pmatrix}
 \quad(m\geq0).
\]
取 $m=10$ 得 $A^{10}=\begin{pmatrix}1&0\\-341&1024\end{pmatrix}$。对角化把矩阵乘方转化为标量乘方，且 $m=0$ 时公式也给出 $I_2$。
\end{solve}


\subsection{习题 6.2 解答}
\label{subsec:xie-exercises-62}

\begin{example}
\label{ex:xie-62-01}
教材习题 6.2 第 1 题：判断例~\ref{ex:xie-61-01} 中的三个矩阵能否相似于对角矩阵，并说明理由。
\end{example}
\begin{solve}
三个矩阵的唯一特征值及其代数重数分别为 $1$（三重）、$-1$（三重）、$0$（三重），而相应特征子空间维数分别为 $2,1,1$。每个矩阵都有一个特征值的几何重数小于代数重数，因此三个矩阵都不可对角化。也可用只有一个特征值的矩阵可对角化时必为纯量阵来判定。
\end{solve}

\begin{example}
\label{ex:xie-62-02}
教材习题 6.2 第 2 题：对下列各矩阵求可逆矩阵 $P_i$，使 $P_i^{-1}A_iP_i$ 为对角矩阵：
\[
 A_1=\begin{pmatrix}-1&4&-2\\-3&4&0\\-3&1&3\end{pmatrix},\quad
 A_2=\begin{pmatrix}1&-2&2\\2&-4&4\\1&-2&2\end{pmatrix},\quad
 A_3=\begin{pmatrix}0&1&0\\1&0&0\\0&0&1\end{pmatrix}.
\]
\end{example}
\begin{solve}
$\chi_{A_1}(t)=(t-1)(t-2)(t-3)$，对应特征子空间分别由 $(1,1,1)^{\mathsf T}$、$(2,3,3)^{\mathsf T}$、$(1,3,4)^{\mathsf T}$ 张成，故可取
\[
 P_1=\begin{pmatrix}1&2&1\\1&3&3\\1&3&4\end{pmatrix},
 \qquad P_1^{-1}A_1P_1=\operatorname{diag}(1,2,3).
\]
$\chi_{A_2}(t)=t^2(t+1)$，其中 $E_0(A_2)$ 由 $(2,1,0)^{\mathsf T}$、$(-2,0,1)^{\mathsf T}$ 张成，$E_{-1}(A_2)$ 由 $(1,2,1)^{\mathsf T}$ 张成，故可取
\[
 P_2=\begin{pmatrix}2&-2&1\\1&0&2\\0&1&1\end{pmatrix},
 \qquad P_2^{-1}A_2P_2=\operatorname{diag}(0,0,-1).
\]
$A_3$ 交换前两个坐标并保持第三个坐标，故取
\[
 P_3=\begin{pmatrix}1&0&1\\1&0&-1\\0&1&0\end{pmatrix},
 \qquad P_3^{-1}A_3P_3=\operatorname{diag}(1,1,-1).
\]
这些矩阵的列向量各自满足对应特征方程且构成一组基，因而都满足所需关系。
\end{solve}

\begin{example}
\label{ex:xie-62-03}
教材习题 6.2 第 3 题：三阶矩阵 $A$ 的特征值为 $1,1,4$，对应特征向量依次为 $(2,1,0)^{\mathsf T}$、$(-1,0,1)^{\mathsf T}$、$(0,1,1)^{\mathsf T}$，求 $A$。
\end{example}
\begin{solve}
将给定向量依次作为列组成 $P=\begin{pmatrix}2&-1&0\\1&0&1\\0&1&1\end{pmatrix}$，其行列式为 $-1$，所以这些向量确为一组基。由 $AP=P\operatorname{diag}(1,1,4)$，得
\[
 A=P\operatorname{diag}(1,1,4)P^{-1}
 =\begin{pmatrix}1&0&0\\-3&7&-3\\-3&6&-2\end{pmatrix}.
\]
给定一组基上的像唯一决定线性变换，所以这个答案唯一。
\end{solve}

\begin{example}
\label{ex:xie-62-04}
教材习题 6.2 第 4 题：设 $Av_j=jv_j$（$j=1,2,3$），其中 $v_1=(1,2,1)^{\mathsf T}$、$v_2=(0,-2,1)^{\mathsf T}$、$v_3=(1,1,2)^{\mathsf T}$，求 $A$。
\end{example}
\begin{solve}
不同特征值对应的这三个非零向量线性无关。令 $P=(v_1,v_2,v_3)$，则
\[
 A=P\operatorname{diag}(1,2,3)P^{-1}
 =\begin{pmatrix}-7&2&4\\-14&5&6\\-13&3&8\end{pmatrix}.
\]
逐列乘回可得 $Av_1=v_1$、$Av_2=2v_2$、$Av_3=3v_3$，也就验证了全部给定条件。
\end{solve}

\begin{example}
\label{ex:xie-62-05}
教材习题 6.2 第 5 题：已知
\[
 A=\begin{pmatrix}1&-1&1\\2&x&-2\\-3&-3&y\end{pmatrix}
 \quad\text{与}\quad B=\operatorname{diag}(2,2,z)
\]
相似，求 $x,y,z$ 以及满足 $P^{-1}AP=B$ 的一个可逆矩阵 $P$。
\end{example}
\begin{solve}
若 $z=2$，则 $B=2I_3$，与它相似的矩阵只能是 $2I_3$，这与 $A$ 的非零非对角元素矛盾。因此 $z\ne2$，相似性给出 $\dim E_2(A)=2$，即 $\operatorname{rank}(A-2I_3)=1$。矩阵
\[
 A-2I_3=\begin{pmatrix}-1&-1&1\\2&x-2&-2\\-3&-3&y-2\end{pmatrix}
\]
的第二行必须等于第一行的 $-2$ 倍，第三行必须等于第一行的 $3$ 倍，故 $x=4,y=5$。比较迹得 $1+x+y=4+z$，所以 $z=6$。此时
\[
 E_2(A)=\operatorname{span}\{(1,0,1)^{\mathsf T},(0,1,1)^{\mathsf T}\},\qquad
 E_6(A)=\operatorname{span}\{(1,-2,3)^{\mathsf T}\}.
\]
因此 $P=\begin{pmatrix}1&0&1\\0&1&-2\\1&1&3\end{pmatrix}$ 的行列式为 $4$，且 $P^{-1}AP=\operatorname{diag}(2,2,6)=B$，也证明这些参数确实充分。
\end{solve}

\begin{example}
\label{ex:xie-62-06}
教材习题 6.2 第 6 题：对 $A=\begin{pmatrix}3&2&-2\\-k&-1&k\\4&2&-3\end{pmatrix}$，求它可对角化的参数条件，并在可对角化时求 $P$ 和对应对角矩阵。
\end{example}
\begin{solve}
直接展开得 $\chi_A(t)=(t-1)(t+1)^2$，与 $k$ 无关。$1$ 为单根，只需检查 $-1$ 的几何重数，而
\[
 A+I_3=\begin{pmatrix}4&2&-2\\-k&0&k\\4&2&-2\end{pmatrix}.
\]
它的第一行非零，第三行与第一行相同；秩为 $1$ 当且仅当第二行是第一行的倍数。比较第二个分量知此倍数只能为零，故恰在 $k=0$ 时成立。若 $k\ne0$，秩为 $2$，几何重数只有 $1$，不可对角化。$k=0$ 时可取
\[
 P=\begin{pmatrix}1&0&1\\-2&1&0\\0&1&1\end{pmatrix},
 \qquad P^{-1}AP=\operatorname{diag}(-1,-1,1),
\]
因为前两列分别为 $-1$ 的两个线性无关特征向量，第三列为 $1$ 的特征向量。
\end{solve}

\begin{example}
\label{ex:xie-62-07}
教材习题 6.2 第 7 题：求矩阵 $A=\begin{pmatrix}1&-1&1\\2&4&-2\\-3&-3&5\end{pmatrix}$ 的 $m$ 次幂，其中 $m\geq0$ 为整数。
\end{example}
\begin{solve}
它就是例~\ref{ex:xie-62-05} 中参数确定后的矩阵，特征值为 $2,6$。多项式
\[
 p_m(t)=2^m+\frac{6^m-2^m}{4}(t-2)
\]
满足 $p_m(2)=2^m$、$p_m(6)=6^m$。在特征向量基下比较对角元，得到
\begin{equation}
\label{eq:xie-exercise-matrix-power}
 A^m=2^mI_3+\frac{6^m-2^m}{4}
 \begin{pmatrix}-1&-1&1\\2&2&-2\\-3&-3&3\end{pmatrix}.
\end{equation}
这也等于 $P\operatorname{diag}(2^m,2^m,6^m)P^{-1}$，其中 $P$ 沿用该例。$m=0,1$ 分别得到 $I_3,A$，与定义吻合。
\end{solve}

\begin{example}
\label{ex:xie-62-08}
教材习题 6.2 第 8 题：证明交换分块对角矩阵 $\operatorname{diag}(A_1,\ldots,A_s)$ 中各方形块的顺序，所得矩阵与原矩阵相似。
\end{example}
\begin{solve}
将各块对应的标准基向量作为一个整体，按新块顺序重新排列，得到一个置换矩阵 $P$。它只是重排基，因而可逆；$P^{-1}AP$ 在新基下的第 $i$ 个对角块就是被移到该位置的原矩阵块。各块大小可以不同，但每个块内部的坐标顺序保持不变，所以得到的恰是题述的新分块对角矩阵。
\end{solve}

\begin{example}
\label{ex:xie-62-09}
教材习题 6.2 第 9 题：证明 $AB-BA$ 不可能相似于 $kI_n$，其中 $k\ne0$。
\end{example}
\begin{solve}
由 $\operatorname{tr}(AB)=\sum_{i,j}a_{ij}b_{ji}=\operatorname{tr}(BA)$，得 $\operatorname{tr}(AB-BA)=0$。若它与 $kI_n$ 相似，迹不变性给出 $0=nk$；但 $K$ 是特征为零的数域，$n\geq1$ 且 $k\ne0$，矛盾。
\end{solve}

\begin{example}
\label{ex:xie-62-10}
教材习题 6.2 第 10 题：若实二阶矩阵 $A$ 满足 $\det A<0$，证明它在实数域上可对角化。
\end{example}
\begin{solve}
$\chi_A(t)=t^2-(\operatorname{tr}A)t+\det A$ 的判别式为 $(\operatorname{tr}A)^2-4\det A>0$，因此有两个不同的实根。由不同特征值的特征向量线性无关，$A$ 在 $\mathbb R$ 上可对角化；两个根的乘积为负，还说明其符号相反。
\end{solve}

\begin{example}
\label{ex:xie-62-11}
教材习题 6.2 第 11 题：设 $A,B$ 有相同的 $n$ 个互不相同的特征值，证明存在 $n$ 阶矩阵 $P,Q$ 使 $A=PQ$、$B=QP$。
\end{example}
\begin{solve}
在包含这些特征值的所讨论数域上，两个矩阵均可对角化；按相同顺序排列对角元，得到它们相似。取可逆矩阵 $S$ 使 $B=S^{-1}AS$，令 $P=S$、$Q=S^{-1}A$，则 $PQ=A$、$QP=B$。这个构造不要求 $A$ 可逆，故允许特征值中包含零。
\end{solve}

\begin{example}
\label{ex:xie-62-12}
教材习题 6.2 第 12 题：证明对合矩阵和幂等矩阵都可对角化。
\end{example}
\begin{solve}
若 $A^2=I_n$，则任意 $v$ 都有分解
\[
 v=\frac{v+Av}{2}+\frac{v-Av}{2},
 \qquad \frac{v+Av}{2}\in E_1(A),\quad
 \frac{v-Av}{2}\in E_{-1}(A).
\]
两个特征子空间交为零，故全空间是它们的直和。若 $A^2=A$，则 $v=Av+(v-Av)$，其中前者属于 $E_1(A)$，后者属于 $E_0(A)$，同样得到直和分解。取这些子空间各自的基即可得到特征向量基。分母 $2$ 可用，是因为数域的特征为零。
\end{solve}

\begin{example}
\label{ex:xie-62-13}
教材习题 6.2 第 13 题：若矩阵 $A,B$ 在 $K$ 上有完全的特征向量系，证明 $\operatorname{diag}(A,B)$ 也有完全的特征向量系。这里同时假设两者的特征多项式在 $K$ 上分裂，即它们在 $K$ 上都有特征向量基。
\end{example}
\begin{solve}
分别取可逆矩阵 $P,Q$ 使 $P^{-1}AP=D_A$、$Q^{-1}BQ=D_B$ 为对角矩阵，则
\[
 \operatorname{diag}(P,Q)^{-1}\operatorname{diag}(A,B)\operatorname{diag}(P,Q)
 =\operatorname{diag}(D_A,D_B),
\]
仍为对角矩阵。因而各块特征向量基补零后合并，就是整个空间的特征向量基；两块有公共特征值也不影响结论。
\end{solve}

\begin{example}
\label{ex:xie-62-14}
教材习题 6.2 第 14 题：设 $A=\begin{pmatrix}I_r&B\\0&-I_{n-r}\end{pmatrix}$，其中 $B$ 为 $r\times(n-r)$ 矩阵，证明 $A$ 可对角化。
\end{example}
\begin{solve}
分块相乘得到 $A^2=I_n$，可直接使用例~\ref{ex:xie-62-12}。也可给出具体矩阵
\[
 P=\begin{pmatrix}I_r&-\tfrac12B\\0&I_{n-r}\end{pmatrix},
 \qquad P^{-1}=\begin{pmatrix}I_r&\tfrac12B\\0&I_{n-r}\end{pmatrix},
\]
乘回得 $P^{-1}AP=\operatorname{diag}(I_r,-I_{n-r})$。若其中一块维数为零，则原矩阵本身已是纯量阵。
\end{solve}

\begin{example}
\label{ex:xie-62-15}
教材习题 6.2 第 15 题：证明只有一个特征值 $\lambda_0$、且全部 $n$ 个复特征值均为 $\lambda_0$ 的非纯量矩阵不可对角化，特别地非零幂零矩阵不可对角化；若实矩阵满足 $A^2+A+I_n=0$，证明它不能在实数域上对角化。
\end{example}
\begin{solve}
若前一种矩阵可对角化，其对角形式只能为 $\lambda_0I_n$，故原矩阵也等于 $\lambda_0I_n$，与非纯量条件矛盾。非零幂零矩阵的特征值全为零，因此是这一结论的特例。对后一种情形，若 $A$ 实可对角化，其每个实对角元 $\lambda$ 都必须满足 $\lambda^2+\lambda+1=0$，而对实数该式等于 $(\lambda+1/2)^2+3/4>0$，矛盾。数域在这里至关重要：在复数域中该二次式有两个不同的根。
\end{solve}

\begin{example}
\label{ex:xie-62-16}
教材习题 6.2 第 16 题：设 $n>1$、$\operatorname{rank}A=1$，证明 $A$ 可对角化，当且仅当 $\operatorname{tr}A\ne0$。
\end{example}
\begin{solve}
秩为 $1$ 保证可写 $A=uv^{\mathsf T}$，其中 $u,v\ne0$。令 $s=v^{\mathsf T}u=\operatorname{tr}A$，则 $A^2=sA$。若 $s\ne0$，$\operatorname{im}A=Ku$，且 $Au=su$，而 $\ker A$ 为 $n-1$ 维；$Ku\cap\ker A=\{0\}$，故 $K^n=Ku\oplus\ker A$ 是特征子空间直和，$A$ 可对角化。若 $s=0$，则 $A^2=0$ 而 $A\ne0$，由上一题不可对角化。
\end{solve}

\section{极小多项式与 Cayley-Hamilton 定理}
\label{sec:xie-minimal-polynomial}
% 依据教材第 283--287 页正文及习题 6.3。

对 $p(t)=c_dt^d+\cdots+c_1t+c_0\in K[t]$，定义 $p(A)=c_dA^d+\cdots+c_1A+c_0I_n$。在同一个矩阵中代入多项式保持加法和乘法，即 $(p+q)(A)=p(A)+q(A)$、$(pq)(A)=p(A)q(A)$；尤其这些矩阵多项式彼此可交换。$M_n(K)$ 的维数为 $n^2$，所以 $I_n,A,\ldots,A^{n^2}$ 必线性相关，说明总存在非零多项式 $p$ 使 $p(A)=0$。本节先找出其中次数最小的基本关系，再证明特征多项式本身也是这样的关系。

\begin{definition}
\label{def:xie-minimal-polynomial}
在所有满足 $p(A)=0$ 的非零多项式中，次数最小的首一多项式称为 $A$ 的\emph{极小多项式}（minimal polynomial），记为 $m_A(t)$，也称最小多项式。对线性变换 $T$，以 $p(T)=0$ 同样定义 $m_T(t)$。
\end{definition}

非零常数多项式的矩阵值为非零常数乘 $I_n$，因此 $\deg m_A\geq1$；初步维数论证给出 $\deg m_A\leq n^2$。极小多项式不一定不可约，“极小”只指所有非零消去多项式中的次数最小。

\begin{theorem}[整除性质与唯一性]
\label{thm:xie-minimal-divisibility}
对任意 $p\in K[t]$，有
\begin{equation}
\label{eq:xie-minimal-divisibility}
 p(A)=0\quad\Longleftrightarrow\quad m_A(t)\mid p(t).
\end{equation}
因此 $m_A$ 唯一。
\end{theorem}
\begin{proof}
若 $m_A\mid p$，代入 $A$ 即得 $p(A)=0$。反之，作带余除法 $p=qm_A+r$，其中 $r=0$ 或 $\deg r<\deg m_A$。代入 $A$ 得 $r(A)=0$，若 $r\ne0$ 就与最小次数矛盾，因此 $r=0$。若 $m$ 和 $\widetilde m$ 都是极小多项式，则彼此整除，故只差非零常数因子；二者首一，所以相等。
\end{proof}

\begin{proposition}[相似与分块对角矩阵]
\label{prop:xie-minimal-block}
相似矩阵具有相同的极小多项式。若 $A=\operatorname{diag}(A_1,\ldots,A_s)$，则
\[
 m_A=\operatorname{lcm}(m_{A_1},\ldots,m_{A_s}),
\]
其中最小公倍式取首一形式。
\end{proposition}
\begin{proof}
若 $B=P^{-1}AP$，则 $p(B)=P^{-1}p(A)P$，所以二者的消去多项式完全相同，从而极小多项式相同。分块对角时 $p(A)=\operatorname{diag}(p(A_1),\ldots,p(A_s))$，它为零等价于每个 $m_{A_i}$ 都整除 $p$，也就等价于这些多项式的首一最小公倍式整除 $p$。
\end{proof}

\begin{example}
\label{ex:xie-minimal-basic}
求纯量阵 $cI_n$、非零二阶幂零矩阵 $N=\begin{pmatrix}0&1\\0&0\end{pmatrix}$ 以及可对角化矩阵的极小多项式。
\end{example}
\begin{solve}
纯量阵满足 $A-cI_n=0$，故 $m_{cI_n}(t)=t-c$。矩阵 $N$ 满足 $N^2=0$ 而 $N\ne0$，其极小多项式整除 $t^2$，又不能是 $t$ 或常数，因此 $m_N(t)=t^2$。若 $A$ 可对角化，全部不同的特征值为 $\lambda_1,\ldots,\lambda_s\in K$，利用相似不变性和分块最小公倍式公式，得到
\begin{equation}
\label{eq:xie-minimal-diagonalizable}
 m_A(t)=\prod_{j=1}^s(t-\lambda_j).
\end{equation}
特征多项式记录各特征值的代数重数，而这里每个不同特征值只出现一次。
\end{solve}

\begin{theorem}[Cayley--Hamilton 定理]
\label{thm:xie-cayley-hamilton}
任意 $A\in M_n(K)$ 都满足自身的特征多项式，即
\begin{equation}
\label{eq:xie-cayley-hamilton}
 \chi_A(A)=0.
\end{equation}
同样，任意有限维线性空间上的线性变换满足 $\chi_T(T)=0$。
\end{theorem}
\begin{proof}
先设 $U$ 为复上三角矩阵，对角元为 $\lambda_1,\ldots,\lambda_n$。令 $F_j=\operatorname{span}\{e_1,\ldots,e_j\}$，其中 $e_i$ 为标准单位列向量，并令 $F_0=\{0\}$。由上三角结构，对每个 $j$ 有
\[
 (U-\lambda_jI_n)F_j\subseteq F_{j-1}.
\]
因此将因子从右向左依次作用，可得
\[
 (U-\lambda_1I_n)\cdots(U-\lambda_nI_n)F_n=\{0\}.
\]
左端的矩阵就是 $\chi_U(U)$，故它是零矩阵。一般的 $A$ 可在 $\mathbb C$ 上三角化，写成 $A=PUP^{-1}$，则 $\chi_A=\chi_U$，且
\[
 \chi_A(A)=P\chi_U(U)P^{-1}=0.
\]
该等式两端原本都在 $M_n(K)$ 中，故也是 $K$ 上的矩阵恒等式。最后，对 $T$ 选一组基，用表示矩阵即得线性变换的结论。
\end{proof}

\begin{corollary}
\label{cor:xie-minimal-characteristic}
$m_A\mid\chi_A$，所以 $1\leq\deg m_A\leq n$；并且 $m_A$ 与 $\chi_A$ 在 $\mathbb C$ 中有相同的根，重数可以不同。
\end{corollary}
\begin{proof}
整除结论由定理~\ref{thm:xie-minimal-divisibility} 和定理~\ref{thm:xie-cayley-hamilton} 得到，故 $m_A$ 的每个根也是 $\chi_A$ 的根。反之，对任一复特征值 $\lambda$ 取非零特征向量 $v$，由 $0=m_A(A)v=m_A(\lambda)v$ 得 $m_A(\lambda)=0$。
\end{proof}

Cayley--Hamilton 定理将最初的 $n^2$ 次上界改进为 $n$ 次，并提供降低矩阵幂次数的方法：对任意 $p$ 作 $p=qm_A+r$、$\deg r<\deg m_A$，便有 $p(A)=r(A)$。这里使用的是多项式代入的代数性质，不能把带有矩阵系数的任意多项式恒等式直接当成标量多项式来代入而忽略乘法顺序。


\subsection{习题 6.3 解答}
\label{subsec:xie-exercises-63}

以下的“公共特征值”在 $\mathbb C$ 中判断；对一般数域 $K$ 上的矩阵，这等价于把特征多项式放到其分裂域中比较根，不能只检查 $K$ 内实际存在的根。

\begin{example}
\label{ex:xie-63-01}
教材习题 6.3 第 1 题：设 $A\in M_n(K)$ 的极小多项式次数为 $d$，证明 $K[A]=\{p(A):p\in K[t]\}$ 是 $M_n(K)$ 的子空间，且 $\dim K[A]=d$。
\end{example}
\begin{solve}
多项式代入保持加法和数乘，所以 $K[A]$ 包含零矩阵并对线性运算封闭。任意 $p$ 除以 $m_A$ 后余式次数小于 $d$，故 $p(A)$ 属于 $\operatorname{span}\{I_n,A,\ldots,A^{d-1}\}$。反之这些幂本来就在 $K[A]$ 中。若它们线性相关，就得到一个次数小于 $d$ 的非零消去多项式，与极小性矛盾。因此它们构成 $K[A]$ 的一组基，维数为 $d$。
\end{solve}

\begin{example}
\label{ex:xie-63-02}
教材习题 6.3 第 2 题：证明矩阵与其转置有相同的极小多项式。
\end{example}
\begin{solve}
对任意 $p\in K[t]$，有 $p(A^{\mathsf T})=p(A)^{\mathsf T}$，因为 $(A^{\mathsf T})^j=(A^j)^{\mathsf T}$，常数项的转置仍为同一个纯量阵。因此 $p(A)=0$ 当且仅当 $p(A^{\mathsf T})=0$，二者消去多项式相同，其唯一首一极小多项式也相同。
\end{solve}

\begin{example}
\label{ex:xie-63-03}
教材习题 6.3 第 3 题：举例说明特征多项式和极小多项式都相同的矩阵未必相似。
\end{example}
\begin{solve}
令 $N=\begin{pmatrix}0&1\\0&0\end{pmatrix}$，取四阶矩阵 $A=\operatorname{diag}(N,N)$、$B=\operatorname{diag}(N,0,0)$，后两个零块为一阶块。二者都是非零矩阵且平方为零，因而极小多项式均为 $t^2$；又都是严格上三角矩阵，特征多项式均为 $t^4$。但 $\operatorname{rank}A=2$、$\operatorname{rank}B=1$，而相似变换不改变秩，所以它们不相似。
\end{solve}

\begin{example}
\label{ex:xie-63-04}
教材习题 6.3 第 4 题：证明 $n$ 阶矩阵的特征多项式满足 $\chi_A(t)\mid m_A(t)^n$。
\end{example}
\begin{solve}
在 $\mathbb C$ 中分解，写 $\chi_A(t)=\prod_i(t-\lambda_i)^{a_i}$，其中 $\lambda_i$ 互异且 $1\leq a_i\leq n$。由推论~\ref{cor:xie-minimal-characteristic}，$m_A$ 有相同的不同根，写成 $m_A(t)=\prod_i(t-\lambda_i)^{b_i}$，其中 $b_i\geq1$。于是 $nb_i\geq n\geq a_i$，逐因式比较得到整除关系。该关系也在 $K[t]$ 中成立：在 $K[t]$ 中作首一多项式除法，其余式在 $\mathbb C[t]$ 中必须为零，故原余式就是零。
\end{solve}

\begin{example}
\label{ex:xie-63-05}
教材习题 6.3 第 5 题：若 $n>1$ 且 $\operatorname{rank}A=1$，证明 $m_A(t)=t^2-(\operatorname{tr}A)t$。
\end{example}
\begin{solve}
写 $A=uv^{\mathsf T}$，其中 $u,v\ne0$，则 $A^2=(v^{\mathsf T}u)A=(\operatorname{tr}A)A$，所以 $q(t)=t^2-(\operatorname{tr}A)t$ 消去 $A$。$A$ 不可能是纯量阵：非零纯量阵的秩是 $n>1$，零阵的秩为零。因此不存在一次首一消去多项式 $t-c$，即 $\deg m_A\geq2$。结合 $m_A\mid q$ 及首一性，得到 $m_A=q$，无论迹是否为零均成立。
\end{solve}

\begin{example}
\label{ex:xie-63-06}
教材习题 6.3 第 6 题：设矩阵 $A,B$ 的特征多项式分别为 $f,g$，极小多项式分别为 $m_A,m_B$，证明下述性质等价：$A,B$ 无公共复特征值；$f,g$ 互素。并证明将后一条件中的任意一个或两个特征多项式换为对应极小多项式，或改为 $f(B),m_A(B),g(A),m_B(A)$ 中任意一个可逆，都得到等价条件。
\end{example}
\begin{solve}
多项式在 $K[t]$ 中互素，当且仅当它们在 $\mathbb C$ 中没有公共根：互素时有 Bézout 恒等式，故无公共根；反之，若有非恒定公因式，它在 $\mathbb C$ 中必有根。$f,m_A$ 的不同根相同，$g,m_B$ 的不同根也相同，因此四种互素条件彼此等价，并均等价于无公共特征值。

由命题~\ref{prop:xie-polynomial-spectrum}，若 $\mu_1,\ldots,\mu_q$ 是 $B$ 的全部复特征值，则 $f(B)$ 的特征值为 $f(\mu_i)$，故 $f(B)$ 可逆当且仅当所有 $f(\mu_i)\ne0$，也就是 $A,B$ 无公共特征值。将 $f$ 换为 $m_A$ 不改变零点集合，交换 $A,B$ 也不改变无公共特征值的条件，因而四种可逆条件也分别等价。
\end{solve}

\begin{example}
\label{ex:xie-63-07}
教材习题 6.3 第 7 题：对任意 $g\in K[t]$，证明 $g(A)$ 可逆，当且仅当 $g$ 与 $m_A$ 互素，也等价于 $g$ 与 $\chi_A$ 互素。
\end{example}
\begin{solve}
若 $g,m_A$ 互素，取 $u,v\in K[t]$ 使 $ug+vm_A=1$，代入 $A$ 得 $u(A)g(A)=I_n$，故 $g(A)$ 可逆，且其逆矩阵就是 $u(A)$。若不互素，则二者有公共复根 $\lambda$；由推论~\ref{cor:xie-minimal-characteristic}，$\lambda$ 是 $A$ 的特征值，取相应非零复特征向量 $x$，便有 $g(A)x=g(\lambda)x=0$，所以 $g(A)$ 不可逆。最后，$m_A$ 与 $\chi_A$ 的不同根相同，故两个互素条件等价。即使 $x$ 不在 $K^n$ 中，存在非零复核也使 $g(A)$ 的行列式为零，从而它在 $K$ 上同样不可逆。
\end{solve}

\begin{example}
\label{ex:xie-63-08}
教材习题 6.3 第 8 题：证明 $A$ 可逆，当且仅当 $m_A$ 的常数项非零。
\end{example}
\begin{solve}
写 $m_A(t)=c_0+tq(t)$。若 $c_0\ne0$，则 $c_0I_n+Aq(A)=0$，所以 $A^{-1}=-q(A)/c_0$。反之，若 $c_0=0$ 且 $A$ 可逆，则 $Aq(A)=0$ 推出 $q(A)=0$，但 $q$ 是比 $m_A$ 次数低的非零多项式，与极小性矛盾。因此 $A$ 可逆时必有 $c_0\ne0$。
\end{solve}

\begin{example}
\label{ex:xie-63-09}
教材习题 6.3 第 9 题：若 $A$ 是 $n$ 阶可逆矩阵，证明存在一个 $n-1$ 次多项式 $g$，使 $A^{-1}=g(A)$。
\end{example}
\begin{solve}
写 $\chi_A(t)=t^n+c_1t^{n-1}+\cdots+c_{n-1}t+c_n$，其中 $c_n=(-1)^n\det A\ne0$。由 Cayley--Hamilton 定理，两边乘 $A^{-1}$ 得
\[
 A^{-1}=-\frac{A^{n-1}+c_1A^{n-2}+\cdots+c_{n-1}I_n}{c_n}.
\]
所以可取 $g(t)=-(t^{n-1}+c_1t^{n-2}+\cdots+c_{n-1})/c_n$，它的最高次项系数为 $-1/c_n\ne0$，次数恰为 $n-1$；$n=1$ 时分子理解为常数 $1$。若改用极小多项式，还可把表示次数降低到小于 $\deg m_A$。
\end{solve}

\begin{example}
\label{ex:xie-63-10}
教材习题 6.3 第 10 题：设 $A\in M_m(K)$、$B\in M_n(K)$ 无公共复特征值，证明方程 $AX=XB$ 只有零解 $X\in M_{m\times n}(K)$。
\end{example}
\begin{solve}
从 $AX=XB$ 归纳得 $A^jX=XB^j$，所以对任意多项式 $p$ 有 $p(A)X=Xp(B)$。取 $p=\chi_A$，利用 Cayley--Hamilton 定理得到 $0=X\chi_A(B)$；由无公共特征值，例~\ref{ex:xie-63-06} 保证 $\chi_A(B)$ 可逆，右乘其逆矩阵便得到 $X=0$。
\end{solve}

\begin{example}
\label{ex:xie-63-11}
教材习题 6.3 第 11 题：在 $M_{m\times n}(K)$ 上定义线性变换 $\mathcal L(X)=AX-XB$。证明 $\mathcal L$ 为线性自同构，当且仅当 $A,B$ 无公共复特征值；此时对每个 $C\in M_{m\times n}(K)$，方程 $AX-XB=C$ 有唯一解。
\end{example}
\begin{solve}
线性性由矩阵乘法的分配律得到。若无公共特征值，上一题给出 $\ker\mathcal L=\{0\}$；它是有限维空间到自身的线性映射，所以单射等价于双射，从而每个 $C$ 有唯一原像。

反之，若有公共复特征值 $\lambda$，取非零复列向量 $u,v$ 使 $Au=\lambda u$、$B^{\mathsf T}v=\lambda v$，其中后者的存在性由 $B$ 与 $B^{\mathsf T}$ 的特征多项式相同保证。则非零矩阵 $X=uv^{\mathsf T}$ 满足 $AX-XB=0$，所以 $\mathcal L$ 在复数域上的表示矩阵不可逆。该表示矩阵的元素原本属于 $K$，它的行列式为零这一事实不随扩域改变，故在 $K$ 上也不可逆。由此得到必要性，而不仅是复数域上的结论。
\end{solve}

\begin{example}
\label{ex:xie-63-12}
教材习题 6.3 第 12 题：设 $T$ 的特征多项式为 $f=f_1f_2$，其中 $f_1,f_2\in K[t]$ 首一且互素，令 $V_i=\ker f_i(T)$。证明 $V_i$ 都是 $T$ 的不变子空间，且
\[
 V=V_1\oplus V_2,\qquad V_1=\operatorname{im}f_2(T),\quad
 V_2=\operatorname{im}f_1(T),\qquad
 \chi_{T|_{V_i}}=f_i\quad(i=1,2).
\]
\end{example}
\begin{solve}
由于 $T$ 与每个 $f_i(T)$ 可交换，若 $f_i(T)v=0$，则 $f_i(T)Tv=T f_i(T)v=0$，故 $V_i$ 不变。取 $a,b\in K[t]$ 使 $af_1+bf_2=1$。Cayley--Hamilton 定理给出 $f_1(T)f_2(T)=0$，所以对任意 $v\in V$，
\[
 v=b(T)f_2(T)v+a(T)f_1(T)v,
\]
其中两项分别属于 $V_1,V_2$。若 $v\in V_1\cap V_2$，对同一 Bézout 恒等式代入便得 $v=0$，因此该和为直和。

由 $f_1(T)f_2(T)=0$ 有 $\operatorname{im}f_2(T)\subseteq V_1$；若 $v\in V_1$，上述恒等式化为 $v=b(T)f_2(T)v=f_2(T)b(T)v$，所以反向包含也成立。交换下标即得另一像空间等式。

最后记 $g_i=\chi_{T|_{V_i}}$。在适合直和分解的基下，表示矩阵分块对角，故 $g_1g_2=f_1f_2$。又 $f_i$ 消去 $T|_{V_i}$，所以 $g_i$ 的每个复根都是 $f_i$ 的根；$f_1,f_2$ 互素保证这两组根不相交。比较乘积中每个根的重数可知 $g_1=f_1,g_2=f_2$，首一性排除了常数因子的差异。若某个 $f_i=1$，则相应 $V_i=\{0\}$，按零维变换的特征多项式为 $1$ 的约定，结论仍成立。
\end{solve}

\section{特征值的估计}
\label{sec:xie-gershgorin}
% 依据教材第 288--291 页正文；按任务范围不录入习题 6.4 和章末复习题。

求出特征多项式的全部根有时很困难，但许多问题只需知道特征值所在的区域，或某个区域内有多少特征值。以下在复数域中讨论。对 $A=(a_{ij})\in M_n(\mathbb C)$，定义第 $i$ 行的非对角元素模长和及相应的闭圆盘为
\begin{equation}
\label{eq:xie-gershgorin-disks}
 R_i=\sum_{j\ne i}|a_{ij}|,
 \qquad D_i=\{z\in\mathbb C:|z-a_{ii}|\leq R_i\}.
\end{equation}
$R_i=0$ 时，$D_i$ 退化为单点。$D_i$ 称为第 $i$ 个\emph{戈尔斯哥林圆盘}（Gershgorin disk），教材简称戈氏圆盘。

\begin{theorem}[Gershgorin 第一圆盘定理]
\label{thm:xie-gershgorin-first}
矩阵 $A$ 的全部特征值都包含在 $\bigcup_{i=1}^nD_i$ 中。
\end{theorem}
\begin{proof}
任取特征值 $\lambda$ 和非零特征向量 $x=(x_1,\ldots,x_n)^{\mathsf T}$，选取 $r$ 使 $|x_r|=\max_i|x_i|>0$。由 $Ax=\lambda x$ 的第 $r$ 行，
\[
 (\lambda-a_{rr})x_r=\sum_{j\ne r}a_{rj}x_j.
\]
运用三角不等式及 $|x_j|\leq|x_r|$，得到
\[
 |\lambda-a_{rr}|\,|x_r|
 \leq\sum_{j\ne r}|a_{rj}|\,|x_j|
 \leq R_r|x_r|.
\]
除以正数 $|x_r|$ 即得 $\lambda\in D_r$。
\end{proof}

最大坐标的下标 $r$ 随特征向量而变，因此该证明保证特征值落入\emph{至少一个}圆盘，不能据此要求它落入每个圆盘。由于 $A$ 与 $A^{\mathsf T}$ 的特征多项式相同，也可以按列定义半径 $C_j=\sum_{i\ne j}|a_{ij}|$，得到另一组包含全部特征值的圆盘；两种圆盘并集的交集仍包含全部特征值。

\begin{example}
\label{ex:xie-gershgorin-four}
估计矩阵
\[
 A=\begin{pmatrix}
 1&0.5&-0.2&-1\\
 0.3&2&-0.2&1.1\\
 -0.5&0.1&-4&0.2\\
 -1&-0.1&0.2&0
 \end{pmatrix}
\]
的特征值范围。
\end{example}
\begin{solve}
按行计算得到
\[
 D_1:\ |z-1|\leq1.7,\qquad D_2:\ |z-2|\leq1.6,
 \qquad D_3:\ |z+4|\leq0.8,\qquad D_4:\ |z|\leq1.3.
\]
因此全部特征值位于这四个圆盘的并集中。$D_1,D_2,D_4$ 彼此连成一个区域，$D_3$ 则与其余圆盘分离；下面的第二圆盘定理将进一步确定这两个区域分别包含三个和一个特征值。
\end{solve}

为了从位置估计推进到计数，需要说明矩阵元素改变时，特征值作为一个整体连续变化。这里的“整体连续”允许重新排列根，不能预先把重根附近的每一个根固定为定义在任意复参数区域上的连续单值函数。

\begin{theorem}[多项式根的整体连续性]
\label{thm:xie-root-continuity}
设 $p(z)=z^n+c_1z^{n-1}+\cdots+c_n$ 的根按重数列为 $\lambda_1,\ldots,\lambda_n$。对每个 $\varepsilon>0$，存在 $\delta>0$，使任何首一多项式 $q(z)=z^n+d_1z^{n-1}+\cdots+d_n$ 满足 $\max_j|d_j-c_j|<\delta$ 时，其根都能排列为 $\mu_1,\ldots,\mu_n$，且 $\max_i|\mu_i-\lambda_i|<\varepsilon$。
\end{theorem}
\begin{proof}
先说明系数有界时所有根一致有界。若 $M=\max_j|d_j|$，则 $|z|>1+M$ 时
\[
 \left|\sum_{j=1}^n d_jz^{n-j}\right|
 \leq M\sum_{j=1}^n|z|^{n-j}<|z|^n,
\]
因而这样的 $z$ 不可能是根。若结论不成立，可取系数趋于 $c_j$ 的多项式序列 $q_k$，使它们的根无论如何排列都不能与 $p$ 的根在固定的 $\varepsilon$ 内配对。任意排列每个 $q_k$ 的根，由根的一致有界性，在 $\mathbb C^n$ 中可取一个收敛子列，其极限记为 $(\nu_1,\ldots,\nu_n)$。对分解式取极限得到
\[
 p(z)=\lim_k q_k(z)=\prod_{i=1}^n(z-\nu_i).
\]
所以 $(\nu_i)$ 只是 $(\lambda_i)$ 的某个排列；沿此子列使用这个固定排列配对，根的距离最终小于 $\varepsilon$，矛盾。
\end{proof}

\begin{theorem}[实区间参数下的连续编号]
\label{thm:xie-eigenvalue-continuous-labeling}
若 $J\subseteq\mathbb R$ 是区间，$A:J\to M_n(\mathbb C)$ 的各矩阵元素均为连续函数，则可选取连续函数 $\lambda_1,\ldots,\lambda_n:J\to\mathbb C$，使每个参数处这 $n$ 个值按重数给出 $A(t)$ 的全部特征值。
\end{theorem}

此处沿用教材将定理~\ref{thm:xie-eigenvalue-continuous-labeling} 作为辅助结论使用，不展开连续编号的证明。它断言\emph{存在一种连续编号}，并非任意编号都连续；参数区域为实区间这一条件也不能随意去掉。例如 $\begin{pmatrix}0&z\\1&0\end{pmatrix}$ 的特征值为 $\pm\sqrt z$，沿复平面中绕原点一周的闭曲线延续一个平方根会变号，因而不能在整个穿孔复平面上作连续单值选择。下面的计数论证实际上只需定理~\ref{thm:xie-root-continuity} 给出的整体连续性。

\begin{theorem}[Gershgorin 第二圆盘定理]
\label{thm:xie-gershgorin-second}
将 $D_1,\ldots,D_n$ 的并集分为连通分支。若一个分支由其中 $k$ 个圆盘组成，则该分支恰好包含 $A$ 的 $k$ 个特征值，按代数重数计数。圆盘按下标计数，即使若干圆盘重合，也分别计入。
\end{theorem}
\begin{proof}
写 $A=D+E$，其中 $D=\operatorname{diag}(a_{11},\ldots,a_{nn})$，$E$ 的对角元为零，并令 $A(t)=D+tE$，$0\leq t\leq1$。它的第 $i$ 个圆盘半径为 $tR_i$，始终包含在固定的 $D_i$ 中，故第一圆盘定理保证所有特征值始终留在原圆盘的并集中。

有限个闭圆盘的并集有有限个连通分支，各分支都是紧集，不同分支间的距离为正。特征多项式系数关于 $t$ 连续，故由根的整体连续性，当 $t$ 的变化充分小时，根可配对得足够接近，不可能从一个分支跳到另一个分支。因此每个分支中的根的数目是 $t$ 的局部常值函数，在连通区间 $[0,1]$ 上为常数。$t=0$ 时矩阵为 $D$，组成该分支的每个圆盘各贡献其圆心处的一个特征值，共 $k$ 个；于是 $t=1$ 时也恰有 $k$ 个。
\end{proof}

两个闭圆盘相交或相切时，都属于同一个连通分支；只有与其余圆盘分离的圆盘组才能单独使用计数结论。回到例~\ref{ex:xie-gershgorin-four}，$D_3$ 中恰有一个特征值，其余三个在 $D_1\cup D_2\cup D_4$ 中。由于 $A$ 是实矩阵，非实特征值成共轭对，而 $D_3$ 关于实轴对称，故该唯一特征值必为实数，且位于区间 $[-4.8,-3.2]$。






